\chapter{WebSocket, REST, TLS et JSON à pleine vitesse}\label{ch:ll:websocket-rest-tls-and-json-at-speed}

Les plateformes crypto et de nombreux courtiers ne parlent pas FIX ni de protocoles binaires avec la plupart de leurs clients~: les données de marché arrivent sous forme de texte JSON
dans des \omterm{def:ll:websocket-rest-tls-and-json-at-speed:ws}{trames WebSocket}, elles-mêmes dans des enregistrements TLS, et les ordres partent sous forme de requêtes HTTP signées ou de JSON sur le même type de connexion.
Un serveur de trading installé dans la région cloud de la plateforme (livre~3, chapitre~26) se trouve à quelques dixièmes de milliseconde de son moteur d'appariement (matching engine).
Sur l'ordinateur portable de ce livre, la référence Python de la construction de ce chapitre consomme près
de \qty{50}{\micro\second} de temps processeur pour recevoir une mise à jour de profondeur et envoyer un ordre (ouverture et fermeture des enregistrements TLS,
démasquage et masquage des trames, décodage du JSON, signature)~; les composants C++ en consomment environ une. La première représente un dixième d'un aller-retour de demi-milliseconde~;
la seconde disparaît dans cet aller-retour. Ce chapitre démonte la pile web couche par couche~: HTTP et le coût de l'ouverture des connexions, TLS et ce que coûtent
sa poignée de main et ses enregistrements, le découpage des \omterm{def:ll:websocket-rest-tls-and-json-at-speed:ws}{trames WebSocket}, le décodage du JSON
à l'aide d'un index structurel, et les signatures de requêtes calculées sans rehacher la clé.

\section{HTTP, REST et réutilisation des connexions}

\begin{definition}[interface REST, connexion persistante, réserve de connexions]\label{def:ll:websocket-rest-tls-and-json-at-speed:rest}
Une \emph{interface REST}\index{interface REST} expose un service sous forme de ressources adressées par des URL et manipulées par des requêtes HTTP
(GET pour lire, POST pour créer, DELETE pour supprimer), chaque requête portant tout ce dont le serveur a besoin pour la traiter. Une \emph{connexion persistante}\index{connexion persistante} transporte
de nombreuses requêtes et réponses HTTP sur une seule connexion TCP (et TLS), au lieu d'en ouvrir une par requête. Une \emph{réserve de connexions}\index{réserve de connexions} (connection pool)
conserve un ensemble de connexions persistantes ouvertes vers un serveur
et en prête une inactive à chaque requête.
\end{definition}

La saisie d'ordres sur une plateforme crypto est généralement une \omterm{def:ll:websocket-rest-tls-and-json-at-speed:rest}{interface REST}~: un ordre est un POST portant les paramètres de l'ordre, la clé d'API du client
et une signature (livre~3, chapitre~15), et la réponse est du JSON. La partie coûteuse n'est pas la requête mais la connexion qui la porte. Ouvrir une connexion coûte
une poignée de main TCP, soit un aller-retour, puis une \omterm{def:ll:websocket-rest-tls-and-json-at-speed:tls}{poignée de main TLS}, soit un autre aller-retour, auxquels s'ajoutent les calculs cryptographiques de l'échange de clés
et de la vérification du certificat~; une \omterm{def:ll:websocket-rest-tls-and-json-at-speed:ws}{WebSocket} ajoute une mise à niveau HTTP, un troisième aller-retour. HTTP/1.1 «~utilise par défaut des \omterm{def:ll:websocket-rest-tls-and-json-at-speed:rest}{connexions persistantes}~» (RFC~9112), et
un client qui les utilise ne paie ces coûts qu'une fois par connexion, et non une fois par ordre.
La réserve de la construction maintient un nombre fixe de connexions ouvertes et fait passer chaque requête d'abord par le régulateur des limites de débit du livre~3,
car les limites d'une plateforme comptent des requêtes et des poids (livre~3, chapitre~15), et non des connexions.

\begin{omfigure}
\begin{tikzpicture}[font=\scriptsize,>=Stealth]
\def\L{0}\def\R{7.5}
\node[font=\footnotesize\bfseries] at (\L,0.45) {client};
\node[font=\footnotesize\bfseries] at (\R,0.45) {plateforme};
\draw[gray,thick] (\L,0.15) -- (\L,-8.1);
\draw[gray,thick] (\R,0.15) -- (\R,-8.1);
\newcommand{\msgR}[3]{\draw[->,#3] (\L,#1) -- node[above,sloped]{#2} (\R,#1-0.3);}
\newcommand{\msgL}[3]{\draw[<-,#3] (\L,#1-0.3) -- node[above,sloped]{#2} (\R,#1);}
\msgR{-0.1}{TCP SYN}{gray}
\msgL{-0.8}{SYN-ACK}{gray}
\msgR{-1.5}{ACK, TLS ClientHello avec une part de clé}{omDef}
\msgL{-2.2}{ServerHello, certificat, Finished}{omDef}
\msgR{-2.9}{Finished, HTTP GET avec Upgrade: websocket}{omMeth}
\msgL{-3.6}{101 Switching Protocols}{omMeth}
\msgL{-4.5}{mise à jour du carnet~: enregistrement TLS [trame WebSocket [JSON]]}{omThm,thick}
\msgR{-5.2}{ordre~: enregistrement TLS [trame masquée [JSON, signature]]}{omThm,thick}
\msgL{-6.4}{ping (toutes les 20~s sur une plateforme)}{gray}
\msgR{-7.1}{pong}{gray}
\draw[decorate,decoration={brace,amplitude=4pt,mirror},thick] (-0.25,-0.1) -- (-0.25,-3.9)
  node[midway,left=5pt,align=right]{trois allers-retours\\ avant le premier\\ message};
\draw[decorate,decoration={brace,amplitude=4pt},thick] (7.75,-4.5) -- (7.75,-5.5)
  node[midway,right=5pt,align=left]{sur une connexion ouverte~:\\ un enregistrement\\ dans chaque sens};
\end{tikzpicture}

\omcaption{Une connexion \omterm{def:ll:websocket-rest-tls-and-json-at-speed:ws}{WebSocket} sur TLS~1.3~: un aller-retour pour TCP, un pour TLS, un pour la mise à niveau HTTP~; ensuite, chaque
message est une \omterm{def:ll:websocket-rest-tls-and-json-at-speed:ws}{trame WebSocket} à l'intérieur d'un enregistrement TLS. Un ordre REST sur une nouvelle connexion paie les deux premiers allers-retours~; sur
une \omterm{def:ll:websocket-rest-tls-and-json-at-speed:rest}{connexion persistante}, aucun.}\label{fig:ll:websocket-rest-tls-and-json-at-speed:setup}
\end{omfigure}

\section{TLS~: la poignée de main et la couche d'enregistrement}

\begin{definition}[poignée de main TLS, reprise de session]\label{def:ll:websocket-rest-tls-and-json-at-speed:tls}
La \emph{poignée de main TLS}\index{poignée de main TLS} est l'échange qui ouvre une connexion TLS~: les deux parties s'accordent sur une version du protocole
et une suite de chiffrement, échangent des parts de clés, le serveur prouve son identité par son certificat, et les deux côtés dérivent les clés qui
protégeront les enregistrements suivants~; en TLS~1.3, elle ne demande qu'un aller-retour. La \emph{reprise de session}\index{reprise de session}
ouvre une nouvelle connexion avec une clé dérivée d'une connexion antérieure (une clé pré-partagée envoyée au client sous forme de ticket), en sautant
le certificat et une partie de l'échange de clés.
\end{definition}

TLS~1.3 (RFC~8446) a réduit la poignée de main complète à un aller-retour et a remplacé les anciens identifiants et tickets de session par
une reprise utilisant une clé pré-partagée~: «~la clé dérivée de la poignée de main initiale sert à amorcer l'état cryptographique
au lieu d'une poignée de main complète~». Un mode sans aller-retour permet à un client qui reprend la session d'envoyer des données avec son premier message, «~au
prix de certaines propriétés de sécurité~» (les données peuvent être rejouées), raison pour laquelle un ordre ne doit jamais emprunter ce mode. Après la poignée de main,
la couche d'enregistrement découpe le flux en enregistrements d'au plus $2^{14}$ octets et protège chacun
avec un chiffrement authentifié.

Le \Cref{tab:ll:websocket-rest-tls-and-json-at-speed:tls} donne les coûts sur la boucle locale de l'ordinateur portable, où le réseau ne coûte
presque rien et où il ne reste que le traitement. Une \omterm{def:ll:websocket-rest-tls-and-json-at-speed:tls}{poignée de main TLS}~1.3 complète avec un certificat à courbes elliptiques prend environ
\qty{2.1}{\milli\second} avec le module \texttt{ssl} de Python (les deux extrémités sur la même machine), une poignée reprise environ
\qty{1.6}{\milli\second}. Un aller-retour de 300 octets sur une connexion ouverte prend environ \qty{45}{\micro\second} en TLS et
\qty{33}{\micro\second} en TCP simple, les deux fois en Python~; la protection elle-même, un chiffrement ou un déchiffrement AES-128-GCM d'un enregistrement de 300 octets
avec OpenSSL appelé depuis C++, reste inférieure à \qty{0.2}{\micro\second}. La poignée de main coûte dix mille fois plus cher qu'un enregistrement~:
c'est la raison de garder les connexions ouvertes.

\begin{table}[ht]
\centering\small
\begin{tabular}{lr}
\toprule
Opération (boucle locale, ordinateur portable) & Coût \\
\midrule
\omterm{def:ll:websocket-rest-tls-and-json-at-speed:tls}{Poignée de main TLS}~1.3 complète (Python \texttt{ssl}, les deux extrémités) & \qty{2.1}{\milli\second} \\
\omterm{def:ll:websocket-rest-tls-and-json-at-speed:tls}{Poignée de main TLS}~1.3 reprise (Python \texttt{ssl}, les deux extrémités) & \qty{1.6}{\milli\second} \\
Aller-retour de 300 octets sur TLS (Python) & \qty{45}{\micro\second} \\
Aller-retour de 300 octets sur TCP simple (Python) & \qty{33}{\micro\second} \\
Chiffrement AES-128-GCM d'un enregistrement de 300 octets (OpenSSL, C++) & \qty{0.16}{\micro\second} \\
Déchiffrement AES-128-GCM d'un enregistrement de 300 octets (OpenSSL, C++) & \qty{0.15}{\micro\second} \\
\bottomrule
\end{tabular}
\caption{TLS sur la boucle locale de l'ordinateur portable~: médianes de 300 poignées de main dans chaque sens et de 5\,000 aller-retours, et de 21 passes de
1\,000 enregistrements. Mesures sur un ordinateur portable (Intel Core Ultra 7 155H) sous WSL2, sans cœurs isolés~; OpenSSL~3.0.2, Python~3.10.
Données~: \texttt{bench\_web.py}, \texttt{measured\_tls.csv} et \texttt{measured\_web.csv}.}\label{tab:ll:websocket-rest-tls-and-json-at-speed:tls}
\end{table}

\section{Découpage des trames WebSocket}

\begin{definition}[WebSocket, trame WebSocket]\label{def:ll:websocket-rest-tls-and-json-at-speed:ws}
Une \emph{WebSocket}\index{WebSocket} est un canal de messages bidirectionnel sur une seule connexion TCP (généralement à l'intérieur de TLS), ouvert par
une requête HTTP qui demande la mise à niveau de la connexion (RFC~6455). Une \emph{trame WebSocket}\index{trame WebSocket} est son unité
sur le fil~: un en-tête de deux octets (un bit de fin de fragment, un code d'opération pour texte, binaire, continuation, fermeture, ping ou pong, un bit de masquage
et une longueur de 7 bits), une longueur étendue de 2 ou 8 octets lorsque la charge utile dépasse 125 octets, une clé de masquage de 4 octets sur les trames
émises par le client, et la charge utile.
\end{definition}

Trois règles de la RFC~6455 façonnent un client. «~Un client DOIT masquer toutes les trames qu'il envoie au serveur~», avec une nouvelle clé aléatoire
par trame, même sur TLS~: la charge utile est combinée par XOR à la clé répétée, un rempart contre les proxys interceptants, et non des données. Un
message peut être découpé en fragments (une première trame portant l'opcode, des trames de continuation, la dernière avec son bit de fin), et
les trames de contrôle (ping, pong, close) peuvent arriver entre les fragments, transportent au plus 125~octets et ne sont jamais fragmentées. Et la
plateforme vérifie la connexion avec des pings auxquels le client doit répondre~: une grande plateforme envoie un ping toutes les 20~secondes et coupe
une connexion qui n'a pas répondu dans la minute. Le décodeur de la construction démasque sur place huit octets à la fois.

\omcode{code/firm/wsclient/cpp/firm_wsclient.hpp}{26}{39}{Démasquage d'une charge utile avec la clé de 4~octets répétée dans un
mot de 64~bits.}\label{lst:ll:websocket-rest-tls-and-json-at-speed:unmask}

Le découpage en trames coûte peu en C++~: environ \qty{24}{\nano\second} pour décoder et démasquer une trame de 300~octets, \qty{16}{\nano\second} pour
encoder et masquer une trame. La référence Python de la construction, qui applique le XOR à la charge utile un octet à la fois, prend environ
\qty{16}{\micro\second} pour la même trame~: plusieurs centaines de fois plus, et le poste le plus lourd de son budget.

\section{JSON à grande vitesse}

\begin{definition}[Index structurel]\label{def:ll:websocket-rest-tls-and-json-at-speed:index}
Un \emph{index structurel} d'un texte JSON est le tableau des positions de ses caractères structurels (accolades, crochets,
deux-points, virgules) et guillemets, calculé en une seule passe vectorisée~; un décodeur parcourt alors l'index au lieu des octets, en lisant chaque
valeur directement entre deux positions connues.
\end{definition}

L'analyseur de Langdale et Lemire (chapitre~14) fonctionne en deux étapes~: la première écrit «~l'emplacement de tous les
caractères structurels \dots\ sous forme d'indices entiers dans un tableau séparé~», la seconde parcourt cet index. La construction applique l'idée à un
format de message, la mise à jour de profondeur qu'une grande plateforme documente pour son flux diff-depth~: un horodatage, un symbole, les premier et dernier
identifiants de mise à jour (\texttt{U}, \texttt{u}) que le constructeur de carnet (book builder) du livre~3 vérifie pour détecter les trous, et des tableaux de cours acheteurs et de cours vendeurs sous forme de
chaînes \texttt{[price, quantity]}. La première étape est le \texttt{positions\_any} du chapitre~14, appliquée aux sept caractères
\texttt{\{\}[]:,} et le guillemet double~; la seconde étape parcourt les positions, la forme étant connue à l'avance, lit chaque nombre là où il se trouve, et
convertit les prix en entiers de $10^{-8}$ sans virgule flottante. Tout élément inattendu (un caractère d'échappement, une clé absente) fait
renvoyer faux, et un analyseur général prend le relais.

\omcode{code/firm/wsclient/cpp/firm_wsclient.hpp}{177}{204}{La seconde étape~: parcours de l'index structurel d'une mise
à jour de profondeur.}\label{lst:ll:websocket-rest-tls-and-json-at-speed:walk}

Sur les 1\,000 mises à jour de profondeur du jeu d'essai (environ 286~octets chacune, jusqu'à dix niveaux par côté), le décodeur indexé met environ
\qty{0.39}{\micro\second} par mise à jour~; un analyseur général qui construit un arbre de tableaux associatifs, de vecteurs et de chaînes et convertit les prix
avec \texttt{strtod} environ \qty{2.6}{\micro\second}, sept fois plus~; le module \texttt{json} de Python et la conversion en
entiers, environ \qty{10}{\micro\second}. Les niveaux décodés, divisés par le tick et le lot de l'instrument, sont ceux qu'applique le
\texttt{firm.wsbook} du livre~3.

\section{Signer les requêtes hors du chemin critique}

\begin{definition}[Signature de requête]\label{def:ll:websocket-rest-tls-and-json-at-speed:signature}
Une \emph{signature de requête} est un code d'authentification de message calculé par le client sur les paramètres d'une requête, avec un secret
partagé avec la plateforme (typiquement HMAC-SHA-256 de la chaîne de requête et du corps, transmis en hexadécimal avec la clé d'API), ce qui permet
à la plateforme de vérifier que la requête vient du titulaire de la clé et n'a pas été modifiée~; un horodatage et une fenêtre de validité incluse dans la
charge utile signée empêchent de rejouer d'anciennes requêtes.
\end{definition}

HMAC (RFC~2104) calcule $H\big((K \oplus \mathit{opad}) \,\|\, H((K \oplus \mathit{ipad}) \,\|\, m)\big)$. Les deux clés complétées
forment chacune un bloc de 64~octets, de sorte que l'état de la fonction de hachage après l'absorption de chacune ne dépend que de la clé~; la RFC note que ces
résultats intermédiaires «~peuvent être précalculés une seule fois au moment de la génération de la clé~». Un signataire qui procède ainsi absorbe la
clé une seule fois, recopie les deux états pour chaque requête et ne hache que le message. Le signataire de la construction le fait en C++, en Rust et
en Python, et il est vérifié contre les vecteurs de la RFC~4231 et contre l'exemple détaillé de la documentation d'une grande plateforme, dont
la clé publiée et la charge utile de l'ordre donnent une signature publiée.

\omcode{code/firm/wsclient/cpp/firm_wsclient.hpp}{279}{301}{Les blocs de remplissage de la clé, hachés une fois à la construction~; chaque signature
recopie deux états.}\label{lst:ll:websocket-rest-tls-and-json-at-speed:hmac}

Les coûts enseignent deux leçons. Précalculer les blocs de remplissage économise environ deux cinquièmes avec le SHA-256 portable de la construction (environ
\qty{0.64}{\micro\second} contre \qty{1.1}{\micro\second} pour un ordre de 119~octets), et quatre cinquièmes avec OpenSSL (environ
\qty{0.19}{\micro\second} contre \qty{0.92}{\micro\second} pour son \texttt{HMAC()} en un seul appel). Et la bibliothèque est plusieurs fois
plus rapide qu'une implémentation portable parce qu'elle utilise les instructions SHA du processeur~: la cryptographie est le seul endroit où
le chemin rapide est du code d'un autre, bien appelé.

\begin{dated}{2026-09}{La signature et les flux d'une plateforme}\label{dat:ll:websocket-rest-tls-and-json-at-speed:venue}
La documentation de l'API spot de Binance (GitHub, consultée en septembre 2026) signe chaque requête avec HMAC-SHA-256 de la chaîne de requête
et du corps, avec le secret de la clé d'API, et borne sa validité au moyen d'un horodatage et d'une fenêtre, \texttt{recvWindow}, de 5\,000
millisecondes par défaut et 60\,000 au plus. Ses flux \omterm{def:ll:websocket-rest-tls-and-json-at-speed:ws}{WebSocket} envoient un ping toutes les 20~secondes, coupent une connexion qui n'a pas
répondu dans la minute, et mettent fin à toute connexion au bout de 24~heures~; son flux diff-depth publie des mises à jour toutes les 100 ou 1\,000
millisecondes.
\end{dated}

La \Cref{fig:ll:websocket-rest-tls-and-json-at-speed:budget} rassemble les éléments~: le temps processeur de la réception d'une mise à jour de profondeur
et de l'envoi d'un ordre, étape par étape, pour les composants C++ de la construction et pour sa référence Python.

\begin{omfigure}
\begin{tikzpicture}
\begin{axis}[width=12cm,height=5.6cm,ybar,bar width=9pt,ymode=log,log origin y=infty,ymin=0.01,ymax=100,
  symbolic x coords={tls_open,frame_decode,json,sign,frame_encode,tls_seal},xtick=data,
  xticklabels={{ouverture\\ enreg.\ TLS},{décodage\\ trame},{décodage\\ JSON},{signature\\ ordre},{encodage,\\ masquage},{scellement\\ enreg.\ TLS}},
  xticklabel style={align=center,font=\scriptsize},enlarge x limits=0.1,
  ylabel={\si{\micro\second} (échelle log)},tick label style={font=\small},label style={font=\small},grid=major,
  grid style={gray!25},legend columns=2,legend style={font=\footnotesize,at={(0.5,-0.3)},anchor=north,draw=none,fill=none,
  /tikz/every even column/.append style={column sep=8pt}},table/col sep=comma]
\addplot[fill=omDef!70,draw=omDef] table[x=step,y=cpp]{figdata/low-latency/17-websocket-rest-tls-and-json-at-speed/measured_budget_steps.csv};
\addplot[fill=omThm!45,draw=omThm] table[x=step,y=python]{figdata/low-latency/17-websocket-rest-tls-and-json-at-speed/measured_budget_steps.csv};
\legend{Composants C++ (0.9~\si{\micro\second} au total),Référence Python (48~\si{\micro\second} au total)}
\end{axis}
\end{tikzpicture}

\omcaption{Le temps processeur d'un ordre \omterm{def:ll:websocket-rest-tls-and-json-at-speed:ws}{WebSocket}~: réception d'une mise à jour de profondeur (enregistrement TLS, trame, JSON) et envoi d'un ordre
(signature, trame masquée, enregistrement TLS). C++~: les décodeurs de la construction et OpenSSL~; Python~: la référence de la construction et \texttt{ssl}
(son coût d'enregistrement estimé au quart du temps d'aller-retour supplémentaire sur TLS). Mesuré sur un ordinateur portable (Intel Core Ultra 7 155H)
sous WSL2, sans cœurs isolés. Données~: \texttt{bench\_web.py}.}\label{fig:ll:websocket-rest-tls-and-json-at-speed:budget}
\end{omfigure}

\section{Tutoriel~: la pile web, mesurée}\label{tut:ll:websocket-rest-tls-and-json-at-speed:tut}

\begin{tutorial}
\textbf{Objectif.} Construire et vérifier chaque couche, puis mesurer le coût d'un ordre \omterm{def:ll:websocket-rest-tls-and-json-at-speed:ws}{WebSocket}. \textbf{État final~:}
le \cref{tab:ll:websocket-rest-tls-and-json-at-speed:tls}, la \cref{fig:ll:websocket-rest-tls-and-json-at-speed:budget}, et
des tests verts en Python, C++ et Rust.
\begin{tutsteps}
\item \textbf{Trames.} Les tests encodent les exemples de la RFC~6455 (un «~Hello~» non masqué et un «~Hello~» masqué, des longueurs nécessitant 16 et 64
bits) et décodent le flux de test écrit par \texttt{make\_ws\_fixtures.py}~: cinquante mises à jour de profondeur, certaines fragmentées en trois
avec un ping entre les fragments, et une fermeture. Python, C++ et Rust délivrent les mêmes messages dans le même ordre.
\item \textbf{Mises à jour de profondeur.} Le décodeur indexé C++ reproduit, niveau par niveau, le décodage de la référence Python
des 1\,000 mises à jour du jeu d'essai.
\item \textbf{Signatures.} Les trois signataires reproduisent les cas de test de la RFC~4231 et l'exemple publié par la plateforme.
\item \textbf{TLS.} \texttt{bench\_web.py} crée un certificat auto-signé avec la commande \texttt{openssl}, lance un serveur TLS~1.3
dans un thread, et mesure les \omterm{def:ll:websocket-rest-tls-and-json-at-speed:tls}{poignées de main TLS} complètes et les \omterm{def:ll:websocket-rest-tls-and-json-at-speed:tls}{reprises de session} (il vérifie que chaque reprise a été acceptée) ainsi que les allers-retours.
\end{tutsteps}
\textbf{Pour aller plus loin.} Démasquer avec des XOR vectoriels de 32~octets (chapitre~14) et mesurer une trame de 16~kilo-octets~; remplacer le
masquage de la référence Python par un XOR au niveau des octets sur des entiers, puis mesurer à nouveau.
\end{tutorial}

\section{Construction~: le client web}\label{bld:ll:websocket-rest-tls-and-json-at-speed:wsclient}

\begin{build}
\bfield{Rôle} La connexion de la firme aux plateformes qui parlent JSON sur \omterm{def:ll:websocket-rest-tls-and-json-at-speed:ws}{WebSocket} et REST~: mises à jour de profondeur décodées pour le
\texttt{firm.wsbook} du livre~3, avec des ordres signés et envoyés par une réserve de \omterm{def:ll:websocket-rest-tls-and-json-at-speed:rest}{connexions persistantes} qui respecte les limites de débit de la plateforme.
\bfield{Interface} Référence Python \texttt{firm\_wsclient}~: \texttt{encode\_frame}, \texttt{decode\_frame},
\texttt{Reassembler}, \texttt{decode\_depth}, \texttt{to\_book\_levels}, \texttt{Signer}, \texttt{Pool}. C++20
\texttt{firm::ws}~: \texttt{encode\_frame}, \texttt{decode\_frame} (démasquage en place), \texttt{Reassembler},
\texttt{decode\_depth} dans un \texttt{Depth} fixe, \texttt{Sha256}, \texttt{HmacSha256}. Rust \texttt{firm\_wsclient}~:
\texttt{encode\_frame}, \texttt{decode\_frame}, \texttt{Sha256}, \texttt{HmacSha256}.
\bfield{Règles} Les trames client sont toujours masquées~; les trames de contrôle ne sont jamais fragmentées et portent au plus 125~octets~; les longueurs sont
codées au minimum~; les prix et les quantités sont des entiers de $10^{-8}$, jamais des flottants~; les pads de la clé sont hachés une seule fois~; la réserve
n'ouvre jamais plus de connexions que sa taille et demande au régulateur de débit avant chaque requête.
\bfield{Tests d'acceptation} \texttt{code/firm/wsclient/}~: exemples de trames de la RFC~6455~; le flux de test reconstitué à l'identique
dans les trois langages~; les 1\,000 mises à jour du jeu d'essai décodées à l'identique en Python et en C++~; les cas de test 1, 2 et 6 de la RFC~4231 et
la signature d'exemple de la plateforme dans les trois langages~; la réutilisation de la réserve, sa capacité et sa limitation de débit.
\bfield{Pour aller plus loin} Démasquage vectorisé~; un lecteur JSON général à index structurel~; la compression deflate par message~; le chemin des ordres à travers
le préréglage crypto du simulateur du livre~10.
\end{build}

\begin{omsources}
\item IETF RFC~6455 (WebSocket), RFC~8446 (TLS~1.3), RFC~2104 (HMAC), RFC~4231 (vecteurs de test HMAC-SHA-256), RFC~9112
(HTTP/1.1).
\item G.~Langdale et D.~Lemire, «~Parsing gigabytes of JSON per second~», \emph{The VLDB Journal} 28(6), 2019.
\item Documentation de l'API spot de Binance (GitHub)~: sécurité des requêtes REST, flux WebSocket.
\end{omsources}

\section{Exercices}

\begin{exercise}[$\star$]\label{exo:ll:websocket-rest-tls-and-json-at-speed:1}
Combien d'octets occupe une trame client portant 300~octets de charge utile~? Et 70\,000~octets~? Et une trame serveur de 100~octets~?
\end{exercise}

\begin{exercise}[$\star$]\label{exo:ll:websocket-rest-tls-and-json-at-speed:2}
Un ordre REST sur une nouvelle connexion doit payer la mise en place de TCP et de TLS. Avec un aller-retour de \qty{0.5}{\milli\second} vers la plateforme et la
poignée de main mesurée, de combien arrive-t-il plus tard que sur une \omterm{def:ll:websocket-rest-tls-and-json-at-speed:rest}{connexion persistante}~?
\end{exercise}

\begin{exercise}[$\star$]\label{exo:ll:websocket-rest-tls-and-json-at-speed:3}
Pourquoi un client ne doit-il jamais envoyer un ordre dans les données à zéro aller-retour (0-RTT) de TLS~1.3~?
\end{exercise}

\begin{exercise}[$\star\star$]\label{exo:ll:websocket-rest-tls-and-json-at-speed:4}
Pourquoi les deux clés complétées de HMAC peuvent-elles être hachées à l'avance, et combien de compressions SHA-256 faut-il pour signer une charge utile de 119~octets
avec et sans ce précalcul~?
\end{exercise}

\begin{exercise}[$\star\star$]\label{exo:ll:websocket-rest-tls-and-json-at-speed:5}
Le prix d'une mise à jour de profondeur est la chaîne \texttt{64123.2}. Que renvoie le décodeur indexé avec une précision de $10^{-8}$, et que peut-il mal tourner si un
décodeur le convertit avec \texttt{strtod} et le multiplie par $10^8$~?
\end{exercise}

\begin{exercise}[$\star\star$]\label{exo:ll:websocket-rest-tls-and-json-at-speed:6}
Un ping arrive entre le premier et le deuxième fragment d'une mise à jour de profondeur. Que doit faire le client, et dans quel ordre~?
\end{exercise}

\begin{exercise}[$\star\star\star$]\label{exo:ll:websocket-rest-tls-and-json-at-speed:7}
\emph{Programmation.} Remplacer le masquage octet par octet de la référence Python par un XOR de deux grands entiers construits avec
\texttt{int.from\_bytes}, vérifier les tests sur le jeu d'essai, et mesurer à nouveau le coût des trames.
\end{exercise}

\begin{exercise}[$\star\star\star$]\label{exo:ll:websocket-rest-tls-and-json-at-speed:8}
\emph{Trouver la faille.} «~Chaque thread de stratégie ouvre sa propre connexion HTTPS pour chaque ordre, de sorte qu'aucun couple de threads
ne se dispute jamais une connexion.~»
\end{exercise}

\section{Problème~: du cloud à la plateforme}

\begin{problem}[{Problème du week-end --- où passe le temps d'un ordre WebSocket}]\label{pb:ll:websocket-rest-tls-and-json-at-speed:1}
Une stratégie tourne dans la même région cloud qu'une plateforme crypto. Supposer un saut réseau à sens unique de \qty{250}{\micro\second} et une
décision de coût négligeable~; prendre les coûts mesurés de la \cref{fig:ll:websocket-rest-tls-and-json-at-speed:budget} et le
\cref{tab:ll:websocket-rest-tls-and-json-at-speed:tls} (\texttt{measured\_budget.csv}, \texttt{measured\_tls.csv}).

\medskip\noindent\textbf{Partie~I --- Le budget.}
\begin{enumerate}
\item Quel est le temps processeur total d'un ordre avec les composants C++, et avec la référence Python~?
\item Quelle étape domine dans chaque cas~?
\item Quelle part du trajet plateforme, hôte, plateforme représente le traitement de l'hôte, dans chaque cas~?
\item Quelle serait la part de Python si seul le découpage en trames était en C++ (le reste inchangé)~?
\end{enumerate}

\medskip\noindent\textbf{Partie~II --- Connexions.}
\begin{enumerate}[resume]
\item Que rajoute l'ouverture d'une nouvelle connexion pour chaque ordre REST, avec une poignée de main complète, et avec une \omterm{def:ll:websocket-rest-tls-and-json-at-speed:tls}{reprise de session}~?
\item Que rajoute une nouvelle connexion \omterm{def:ll:websocket-rest-tls-and-json-at-speed:ws}{WebSocket} avant son premier message~?
\item Combien d'ordres par seconde une connexion pourrait-elle transporter si chacun ouvrait une nouvelle connexion~?
\item Qu'implique la limite de 24~heures sur les connexions de la plateforme pour une stratégie qui tourne pendant des jours~?
\end{enumerate}

\medskip\noindent\textbf{Partie~III --- Signatures et données.}
\begin{enumerate}[resume]
\item Pourquoi la signature ne peut-elle pas être calculée avant la décision, et qu'est-ce qui peut l'être~?
\item À 100 mises à jour de profondeur par seconde et par instrument sur 50 instruments, quelle part d'un cœur le décodage prend-il avec chacun des
décodeurs~?
\item Que permet de détecter le client grâce à la paire d'identifiants de mise à jour \texttt{(U, u)}, et que fait-il alors~?
\item Pourquoi l'horodatage est-il inclus dans la charge utile signée~?
\end{enumerate}

\medskip\noindent\textbf{Partie~IV --- Le verdict.}
\begin{enumerate}[resume]
\item Énoncer le \emph{résultat nommé}~: la part du tick-to-trade consacrée au déchiffrement, au découpage en trames, au JSON et à la signature, face au
saut réseau, et ce que la réutilisation des connexions économise par ordre.
\item Quand la référence Python suffit-elle~?
\item Où placerait-on le travail TLS si le budget était encore plus serré~?
\item Que mesurerait-on sur le serveur réel avant de faire confiance à ces chiffres obtenus sur ordinateur portable~?
\item Pourquoi le saut réseau n'apparaît-il pas dans le budget processeur, et pourquoi reste-t-il important~?
\item Quelle partie de cette pile supprime le protocole binaire d'une plateforme~?
\item Comment tester le comportement de la réserve lorsque la plateforme répond 429~?
\item En une phrase~: qu'est-ce qui est rapide dans une pile web rapide~?
\end{enumerate}
\end{problem}

\section{Questions d'entretien}

\begin{interviewq}[$\star$ \iqroles{développeur}]\label{iq:ll:websocket-rest-tls-and-json-at-speed:1}
Pourquoi réutiliser les connexions HTTP pour la saisie des ordres~? Que coûte une nouvelle connexion~?
\end{interviewq}

\begin{interviewq}[$\star\star$ \iqroles{développeur}]\label{iq:ll:websocket-rest-tls-and-json-at-speed:2}
Décrivez une \omterm{def:ll:websocket-rest-tls-and-json-at-speed:ws}{trame WebSocket}. Pourquoi les trames client sont-elles masquées, et que coûte le masquage~?
\end{interviewq}

\begin{interviewq}[$\star\star$ \iqroles{développeur}]\label{iq:ll:websocket-rest-tls-and-json-at-speed:3}
Comment analyseriez-vous les mises à jour de profondeur JSON de la plateforme, aussi vite que possible~?
\end{interviewq}

\begin{interviewq}[$\star\star$ \iqroles{développeur}]\label{iq:ll:websocket-rest-tls-and-json-at-speed:4}
Comment calcule-t-on une signature HMAC de requête, et comment rendriez-vous la signature peu coûteuse~?
\end{interviewq}

\begin{interviewq}[$\star\star$ \iqroles{développeur}]\label{iq:ll:websocket-rest-tls-and-json-at-speed:5}
Qu'économise la \omterm{def:ll:websocket-rest-tls-and-json-at-speed:tls}{reprise de session} TLS~1.3, et que risquent les données à zéro aller-retour (0-RTT)~?
\end{interviewq}

\begin{interviewq}[$\star\star\star$ \iqroles{développeur, chercheur}]\label{iq:ll:websocket-rest-tls-and-json-at-speed:6}
Les ordres de votre stratégie crypto arrivent 3~ms après les données de marché qui les ont déclenchés, alors que la plateforme n'est qu'à 0.3~ms. Où
regardez-vous~?
\end{interviewq}
